% Options for packages loaded elsewhere
\PassOptionsToPackage{unicode}{hyperref}
\PassOptionsToPackage{hyphens}{url}
\PassOptionsToPackage{dvipsnames,svgnames,x11names}{xcolor}
%
\documentclass[
  10pt,
  a4paper,
]{article}
\usepackage{amsmath,amssymb}
\usepackage{iftex}
\ifPDFTeX
  \usepackage[T1]{fontenc}
  \usepackage[utf8]{inputenc}
  \usepackage{textcomp} % provide euro and other symbols
\else % if luatex or xetex
  \usepackage{unicode-math} % this also loads fontspec
  \defaultfontfeatures{Scale=MatchLowercase}
  \defaultfontfeatures[\rmfamily]{Ligatures=TeX,Scale=1}
\fi
\usepackage{lmodern}
\ifPDFTeX\else
  % xetex/luatex font selection
\fi
% Use upquote if available, for straight quotes in verbatim environments
\IfFileExists{upquote.sty}{\usepackage{upquote}}{}
\IfFileExists{microtype.sty}{% use microtype if available
  \usepackage[]{microtype}
  \UseMicrotypeSet[protrusion]{basicmath} % disable protrusion for tt fonts
}{}
\makeatletter
\@ifundefined{KOMAClassName}{% if non-KOMA class
  \IfFileExists{parskip.sty}{%
    \usepackage{parskip}
  }{% else
    \setlength{\parindent}{0pt}
    \setlength{\parskip}{6pt plus 2pt minus 1pt}}
}{% if KOMA class
  \KOMAoptions{parskip=half}}
\makeatother
\usepackage{xcolor}
\usepackage[margin=21mm]{geometry}
\usepackage{color}
\usepackage{fancyvrb}
\newcommand{\VerbBar}{|}
\newcommand{\VERB}{\Verb[commandchars=\\\{\}]}
\DefineVerbatimEnvironment{Highlighting}{Verbatim}{commandchars=\\\{\}}
% Add ',fontsize=\small' for more characters per line
\newenvironment{Shaded}{}{}
\newcommand{\AlertTok}[1]{\textcolor[rgb]{1.00,0.00,0.00}{\textbf{#1}}}
\newcommand{\AnnotationTok}[1]{\textcolor[rgb]{0.38,0.63,0.69}{\textbf{\textit{#1}}}}
\newcommand{\AttributeTok}[1]{\textcolor[rgb]{0.49,0.56,0.16}{#1}}
\newcommand{\BaseNTok}[1]{\textcolor[rgb]{0.25,0.63,0.44}{#1}}
\newcommand{\BuiltInTok}[1]{\textcolor[rgb]{0.00,0.50,0.00}{#1}}
\newcommand{\CharTok}[1]{\textcolor[rgb]{0.25,0.44,0.63}{#1}}
\newcommand{\CommentTok}[1]{\textcolor[rgb]{0.38,0.63,0.69}{\textit{#1}}}
\newcommand{\CommentVarTok}[1]{\textcolor[rgb]{0.38,0.63,0.69}{\textbf{\textit{#1}}}}
\newcommand{\ConstantTok}[1]{\textcolor[rgb]{0.53,0.00,0.00}{#1}}
\newcommand{\ControlFlowTok}[1]{\textcolor[rgb]{0.00,0.44,0.13}{\textbf{#1}}}
\newcommand{\DataTypeTok}[1]{\textcolor[rgb]{0.56,0.13,0.00}{#1}}
\newcommand{\DecValTok}[1]{\textcolor[rgb]{0.25,0.63,0.44}{#1}}
\newcommand{\DocumentationTok}[1]{\textcolor[rgb]{0.73,0.13,0.13}{\textit{#1}}}
\newcommand{\ErrorTok}[1]{\textcolor[rgb]{1.00,0.00,0.00}{\textbf{#1}}}
\newcommand{\ExtensionTok}[1]{#1}
\newcommand{\FloatTok}[1]{\textcolor[rgb]{0.25,0.63,0.44}{#1}}
\newcommand{\FunctionTok}[1]{\textcolor[rgb]{0.02,0.16,0.49}{#1}}
\newcommand{\ImportTok}[1]{\textcolor[rgb]{0.00,0.50,0.00}{\textbf{#1}}}
\newcommand{\InformationTok}[1]{\textcolor[rgb]{0.38,0.63,0.69}{\textbf{\textit{#1}}}}
\newcommand{\KeywordTok}[1]{\textcolor[rgb]{0.00,0.44,0.13}{\textbf{#1}}}
\newcommand{\NormalTok}[1]{#1}
\newcommand{\OperatorTok}[1]{\textcolor[rgb]{0.40,0.40,0.40}{#1}}
\newcommand{\OtherTok}[1]{\textcolor[rgb]{0.00,0.44,0.13}{#1}}
\newcommand{\PreprocessorTok}[1]{\textcolor[rgb]{0.74,0.48,0.00}{#1}}
\newcommand{\RegionMarkerTok}[1]{#1}
\newcommand{\SpecialCharTok}[1]{\textcolor[rgb]{0.25,0.44,0.63}{#1}}
\newcommand{\SpecialStringTok}[1]{\textcolor[rgb]{0.73,0.40,0.53}{#1}}
\newcommand{\StringTok}[1]{\textcolor[rgb]{0.25,0.44,0.63}{#1}}
\newcommand{\VariableTok}[1]{\textcolor[rgb]{0.10,0.09,0.49}{#1}}
\newcommand{\VerbatimStringTok}[1]{\textcolor[rgb]{0.25,0.44,0.63}{#1}}
\newcommand{\WarningTok}[1]{\textcolor[rgb]{0.38,0.63,0.69}{\textbf{\textit{#1}}}}
\usepackage{longtable,booktabs,array}
\usepackage{calc} % for calculating minipage widths
% Correct order of tables after \paragraph or \subparagraph
\usepackage{etoolbox}
\makeatletter
\patchcmd\longtable{\par}{\if@noskipsec\mbox{}\fi\par}{}{}
\makeatother
% Allow footnotes in longtable head/foot
\IfFileExists{footnotehyper.sty}{\usepackage{footnotehyper}}{\usepackage{footnote}}
\makesavenoteenv{longtable}
\setlength{\emergencystretch}{3em} % prevent overfull lines
\providecommand{\tightlist}{%
  \setlength{\itemsep}{0pt}\setlength{\parskip}{0pt}}
\setcounter{secnumdepth}{-\maxdimen} % remove section numbering
\usepackage{amsmath,amssymb,mathtools}
\usepackage{fontspec}
\usepackage{unicode-math}
\setmainfont{Noto Serif}
\setsansfont{Noto Sans}
\setmonofont{DejaVu Sans Mono}[Scale=0.78]
\setmathfont{latinmodern-math.otf}
\usepackage{xurl}
\usepackage{fvextra}
\DefineVerbatimEnvironment{Highlighting}{Verbatim}{breaklines,commandchars=\\\{\}}
\usepackage{microtype}
\usepackage{fancyhdr}
\usepackage{etoolbox}
\usepackage{needspace}
\pagestyle{fancy}
\fancyhf{}
\fancyhead[L]{\footnotesize Research review edition}
\fancyhead[R]{\footnotesize 26 September 2026}
\fancyfoot[R]{\thepage}
\renewcommand{\headrulewidth}{0.3pt}
\setlength{\headheight}{14pt}
\setlength{\emergencystretch}{3em}
\setlength{\parskip}{5pt}
\setlength{\tabcolsep}{4pt}
\renewcommand{\arraystretch}{1.13}
\AtBeginEnvironment{longtable}{\footnotesize}
\AtBeginEnvironment{quote}{\small}
\allowdisplaybreaks
\makeatletter
\renewcommand\section{\@startsection{section}{1}{\z@}{-2.4ex \@plus -1ex \@minus -.2ex}{1.2ex \@plus .2ex}{\normalfont\Large\bfseries}}
\renewcommand\subsection{\@startsection{subsection}{2}{\z@}{-2ex \@plus -.8ex \@minus -.2ex}{.9ex \@plus .2ex}{\normalfont\large\bfseries}}
\makeatother
\ifLuaTeX
  \usepackage{selnolig}  % disable illegal ligatures
\fi
\usepackage{bookmark}
\IfFileExists{xurl.sty}{\usepackage{xurl}}{} % add URL line breaks if available
\urlstyle{same}
\hypersetup{
  colorlinks=true,
  linkcolor={Maroon},
  filecolor={Maroon},
  citecolor={Blue},
  urlcolor={black},
  pdfcreator={LaTeX via pandoc}}

\author{}
\date{}

\begin{document}

\section{Frame-Aligned Records: Supplementary
Information}\label{frame-aligned-records-supplementary-information}

\subsection{Calculations, study definitions and
reproducibility}\label{calculations-study-definitions-and-reproducibility}

\textbf{Supplementary Information v1.2 \textbar{} 26 September 2026}

\textbf{James McGaughran}

\textbf{Scope.} This supplement supports the methods paper. Scientific
calculations and model definitions come first; dated lineage is in the
final provenance section. A reported historical run, a fresh same-code
replay and a separate numerical implementation have different evidential
meanings.

\textbf{Version crosswalk.} Portfolio v1.8 identifies the distribution;
\emph{Frame-Aligned Records} v0.5, this supplement v1.2, the results log
v1.2, protocol reader edition v1.3 and \emph{Absolute Infinite Union}
v1.4 identify its components. CAL-1.6 and similar identifiers name fixed
calculation suites, not alternative current manuscripts. A historical
label means a preserved prior state, even when its date is only one day
earlier. The authoritative executed specification remains protocol v1.2
and its original pre-execution receipt; later reader edits do not modify
that record. Finalization retains these version labels and records
changes through source hashes and a dated change record.

\subsection{S1. Numerical provenance for the methods
paper}\label{s1.-numerical-provenance-for-the-methods-paper}

The source run labels S01 to S06 refer to the preserved six-script
calibration suite under
\nolinkurl{historical/supporting_lineage/v1_6_code/legacy_v1_5/legacy_claude_v0_2_0/}.
The original display values and sampling histories are retained, not all
rerun during this editorial release. The newer calibration and pilot
paths are listed explicitly below.

\begin{longtable}[]{@{}
  >{\raggedright\arraybackslash}p{(\columnwidth - 4\tabcolsep) * \real{0.3333}}
  >{\raggedright\arraybackslash}p{(\columnwidth - 4\tabcolsep) * \real{0.3333}}
  >{\raggedright\arraybackslash}p{(\columnwidth - 4\tabcolsep) * \real{0.3333}}@{}}
\toprule\noalign{}
\begin{minipage}[b]{\linewidth}\raggedright
Statement
\end{minipage} & \begin{minipage}[b]{\linewidth}\raggedright
Value(s)
\end{minipage} & \begin{minipage}[b]{\linewidth}\raggedright
Script and output key
\end{minipage} \\
\midrule\noalign{}
\endhead
\bottomrule\noalign{}
\endlastfoot
Lemma 1 counterexample & \(L=401\) & S02
\nolinkurl{L1_counterexample_published_expression_without_identity} \\
Lemma 3 invariance; generic change & \(\le6\times10^{-17}\); \(0.144\) &
S02 \nolinkurl{L3_*} \\
Proposition 1(c) & \(1.0\) & S02
\nolinkurl{P1_commuting_unitary_changes_pure_ground_state_by} \\
Proposition 1, Remark 2 & \(1.0, 0.399, 0, 0.399, 1.0\) bits & S02
\nolinkurl{P1_stationary_average_conceals_conditional_dynamics} \\
Proposition 2 & \(0.5270653410\); \(0.3400520132\); \(-0.7615941560\);
\(0\); \(1\to0\) bits & S01 \nolinkurl{P2_*} \\
Lemma 4 example & \(0.6428571429\) (formula and enumeration) & S01
\nolinkurl{hypergeometric_*} \\
Proposition 3 checks, Table 1 & as tabulated & S03
\nolinkurl{A_haar_fragments}, \nolinkurl{B_dynamics} \\
Table 2 & as tabulated & S04 \nolinkurl{duplicate_inflation} \\
Hoeffding half-widths & \(0.200\), \(0.100\), \(0.050\) & S01
\nolinkurl{hoeffding_half_width} \\
Supplement Tables S4 to S6 & as tabulated & S05 \nolinkurl{e1},
\nolinkurl{e2}, \nolinkurl{e3}, \nolinkurl{e4} in the exploratory
output \\
Supplement Section S2.5 & as stated & S05 replication output \\
Legacy S06 span checks & \(0\); \(7\times10^{-16}\); \(0.4166\) versus
\(0.4170\) & S06 \\
Mixed twirl, historical v1.5 & 2,000 draws; mean \(0.0557464152\); SE
\(0.0007264476\) & Archived W03 v1.1, Section 6; not the CAL-1.6
sample \\
Mixed twirl, CAL-1.6 & 1,200 draws; mean \(0.05550839355\); SE
\(0.00093479233\); exact \(0.056\) &
\nolinkurl{historical/supporting_lineage/v1_6_data/results.json},
\nolinkurl{mixed_summary}; full draws in \nolinkurl{mixed_draws.csv} \\
Proposition 5a endpoint constant & \(\sqrt{1.0025}=1.0012492197\) at
\(n=2,4,6\) &
\nolinkurl{historical/supporting_lineage/v1_6_data/localized_constants.csv} \\
Proposition 6 spectator & error difference approximately
\(8.44\times10^{-15}\) against a loose bound \(78.53981634\) &
\nolinkurl{historical/supporting_lineage/v1_6_data/results.json},
\nolinkurl{spectator} \\
Unequal certificate example & \(0.0995\) and \(0.104\), same true error
\(0.099\) & Section 7.3.1, direct arithmetic; not two physical
observations \\
Proposition 7, Figure 6 & capped bound \(1\) through \(n=18\); log10
bound approximately \(-56.85\) at \(n=19\) &
\nolinkurl{data/prop7_thresholds.csv}; direct envelope, not state
simulation \\
Source-off feasibility pilot & 192 independent units in two size strata;
separate outcome definitions & \nolinkurl{data/pilot_run/results.json},
\nolinkurl{per_unit.json}; results log, not an LRA test \\
Post hoc sign analysis & primary raw \(p=5.4592\times10^{-5}\); four-row
Holm \(p=1.6377\times10^{-4}\) &
\nolinkurl{data/review_additions/sign_tests.csv};
\nolinkurl{code/analyze_review_additions.py} \\
Range-only precision budget & 32,791 units for half-width 0.015, not a
power claim & \nolinkurl{data/review_additions/review_analysis.json},
\nolinkurl{precision} \\
Initial eligibility and transient access & 5 and 10 eligible site/unit
pairs in native arms briefly exceed 0.8 &
\nolinkurl{data/review_additions/initial_floor_diagnostics.csv} \\
Source encoding in clean transport control & phase code \(d=|f|\);
vacuum/excitation code \(d=|f|^2\) &
\nolinkurl{data/review_additions/transport_encodings.csv}; separate
full-space checks \\
\end{longtable}

\textbf{Additional run crosswalk}

\begin{longtable}[]{@{}
  >{\raggedright\arraybackslash}p{(\columnwidth - 4\tabcolsep) * \real{0.3333}}
  >{\raggedright\arraybackslash}p{(\columnwidth - 4\tabcolsep) * \real{0.3333}}
  >{\raggedright\arraybackslash}p{(\columnwidth - 4\tabcolsep) * \real{0.3333}}@{}}
\toprule\noalign{}
\begin{minipage}[b]{\linewidth}\raggedright
Target
\end{minipage} & \begin{minipage}[b]{\linewidth}\raggedright
Source and data
\end{minipage} & \begin{minipage}[b]{\linewidth}\raggedright
Boundary
\end{minipage} \\
\midrule\noalign{}
\endhead
\bottomrule\noalign{}
\endlastfoot
Factorial four/six/eight-qubit means &
\nolinkurl{data/reviewer_replay/factorial_results.json};
\nolinkurl{factorial_n*.npz} & Same-code reproduction of an exploratory
review extension \\
Larger source-off sample &
\nolinkurl{data/reviewer_replay/reimpl_id1_results.json} & Fresh-stream
reimplementation, not the original sample \\
Third-frame means &
\nolinkurl{data/reviewer_replay/factorial_neutral.json} & Post hoc
comparator, no guaranteed neutrality \\
Paired range and finite-family intervals &
\nolinkurl{data/factorial_analysis/analysis.json};
\nolinkurl{code/analyze_factorial.py} & Derived secondary sampling
bounds \\
Alternative numerical method &
\nolinkurl{verification/factorial_independent_checks.json} & Selected
trajectories, not a second full population study \\
\end{longtable}

\subsection{S2. Historical continuously driven
simulation}\label{s2.-historical-continuously-driven-simulation}

The material below retains the complete exploratory model and outcomes
for reproducibility. It is not another experiment or a retrospective
confirmation of the newer source-off target. References to original
Sections 2 to 5 concern the methods framework; references to Section 6.x
are mapped to S2.x below. Its reported causal interpretations are
descriptive and family-specific, not general theorems.

\subsubsection{S2.0 Status}\label{s2.0-status}

\textbf{Status: EXPLORATORY.} This section reports computed results for
one small model class. The configuration was changed after the first
results were inspected (Section S2.2), so they are not a confirmatory
test of the prospectively stated population conjectures. They can
nevertheless supply counterexamples to a universal claim and guide later
design; exploratory does not mean scientifically worthless. They are
included because they show concretely how the design choices of Sections
2 and 5 decide outcomes. The frozen configuration was then re-run with
fresh seeds (Section S2.5).

The notation below uses \(K_{\mathrm C}\) for the source-coupling
operator, reserving C for an alternate frame label. Preserved scripts
and original outputs still use their original variable C; this is a
notation change only.

\subsubsection{S2.1 Model}\label{s2.1-model}

\begin{itemize}
\tightlist
\item
  \textbf{System and source.} \(n=3\) system qubits and a binary source
  with equal priors; \(H_x=H_S+(-1)^xK_{\mathrm C}\).
\item
  \textbf{Internal Hamiltonian.} \(H_S\) is a random nearest-neighbour
  2-local chain: Gaussian coefficients on all two-body Pauli products on
  adjacent pairs, plus one-body terms at half weight. It is normalised
  to Frobenius norm \(\lambda\).
\item
  \textbf{Coupling and time window.} The source coupling is the star
  operator \(K_{\mathrm C}=\sum_{j=1}^3X_j\), of Frobenius norm
  \(\sqrt{24}\approx4.90\). Times are \(t\in\{0.6,0.7,0.8,0.9,1.0\}\).
\item
  \textbf{Candidates.} There are \(N=20\) candidates: the identity (the
  frame in which \(H_S\) and \(K_{\mathrm C}\) are local), a preparation
  frame \(W\), and 18 Haar-random unitaries.
\item
  \textbf{Locality.} \(L\) is the internal score \(L^{\mathrm{int}}\) on
  the span of Pauli strings of weight at most 2, identity included.
\item
  \textbf{Records (simplified).} The record score is the time-averaged
  number of single qubits whose Helstrom error is at most \(0.1\). This
  is a simplification of the disjoint, persistent redundancy of Section
  2.5.
\item
  \textbf{Near-optimal sets.} \(\mathcal N_L\) requires
  \(L\le\min L+0.02\) and \(L\le0.6\); \(\mathcal N_R\) requires
  \(R=\max R\) and \(R\ge0.5\).
\item
  \textbf{Arms.} In the \emph{shared-frame} arm the initial state is
  \(\left| 000\right\rangle\). In the \emph{decoupled-preparation} arm
  it is \(W\left| 000\right\rangle\) with \(W\) Haar-random, so by
  Section 4.3 it is Haar distributed. The coupling stays in the identity
  frame in both arms.
\end{itemize}

\subsubsection{S2.2 Development history}\label{s2.2-development-history}

The first configuration used a single-qubit coupling
\(X\otimes I\otimes I\), \(\lambda=4\) and nine times in \([0.3,1.5]\).
Records qualified in 1\% of shared-frame models and in none of the
decoupled models (\(A=0.010\) and \(0.000\)). The coupling was then
changed to the star operator, the time window was narrowed, and a grid
over \(\lambda\) was added. All results below use the changed
configuration, which was frozen on 25 September 2026.

\subsubsection{S2.3 Shared versus decoupled
preparation}\label{s2.3-shared-versus-decoupled-preparation}

\Needspace{12\baselineskip}

\textbf{Table S4. Association results by arm (200 models per cell;
Hoeffding 95\% intervals clipped at \(1-1/N=0.95\)).}

\begin{longtable}[]{@{}
  >{\raggedright\arraybackslash}p{(\columnwidth - 12\tabcolsep) * \real{0.1429}}
  >{\raggedright\arraybackslash}p{(\columnwidth - 12\tabcolsep) * \real{0.1429}}
  >{\raggedright\arraybackslash}p{(\columnwidth - 12\tabcolsep) * \real{0.1429}}
  >{\raggedright\arraybackslash}p{(\columnwidth - 12\tabcolsep) * \real{0.1429}}
  >{\raggedright\arraybackslash}p{(\columnwidth - 12\tabcolsep) * \real{0.1429}}
  >{\raggedright\arraybackslash}p{(\columnwidth - 12\tabcolsep) * \real{0.1429}}
  >{\raggedright\arraybackslash}p{(\columnwidth - 12\tabcolsep) * \real{0.1429}}@{}}
\toprule\noalign{}
\begin{minipage}[b]{\linewidth}\raggedright
\(\lambda\)
\end{minipage} & \begin{minipage}[b]{\linewidth}\raggedright
Arm
\end{minipage} & \begin{minipage}[b]{\linewidth}\raggedright
Records qualify
\end{minipage} & \begin{minipage}[b]{\linewidth}\raggedright
\(\bar A\)
\end{minipage} & \begin{minipage}[b]{\linewidth}\raggedright
\(\bar q\)
\end{minipage} & \begin{minipage}[b]{\linewidth}\raggedright
\(\hat d\)
\end{minipage} & \begin{minipage}[b]{\linewidth}\raggedright
95\% interval
\end{minipage} \\
\midrule\noalign{}
\endhead
\bottomrule\noalign{}
\endlastfoot
0.3 & shared & 1.000 & 1.000 & 0.050 & 0.950 & {[}0.758, 0.950{]} \\
0.3 & decoupled & 0.095 & 0.050 & 0.005 & 0.045 & {[}-0.147, 0.237{]} \\
1.0 & shared & 1.000 & 1.000 & 0.050 & 0.950 & {[}0.758, 0.950{]} \\
1.0 & decoupled & 0.060 & 0.025 & 0.003 & 0.022 & {[}-0.170, 0.214{]} \\
4.0 & shared & 0.320 & 0.255 & 0.016 & 0.239 & {[}0.047, 0.431{]} \\
4.0 & decoupled & 0.030 & 0.005 & 0.002 & 0.004 & {[}-0.189, 0.196{]} \\
\end{longtable}

In the shared arm, the most local candidate hosts the records in every
model at \(\lambda\le1\), and the descriptive estimate is \(0.95\) in
these sampled cases. This agreement is facilitated by the shared inputs,
not proved for every model by a coupling bound. At \(\lambda=4\) the
internal dynamics is comparable in norm to the coupling, and records
qualify in only 32\% of models. In the decoupled arm, records qualify
rarely. In the reported cases they preferentially sit in the identity
frame, where the coupling remains local. This association is consistent
with the supplied interaction structure, not an isolated causal
estimate.

\subsubsection{S2.4 Moving the preparation frame and the coupling
frame}\label{s2.4-moving-the-preparation-frame-and-the-coupling-frame}

\Needspace{12\baselineskip}

\textbf{Table S5. Preparation frame rotated by \(W=e^{-i\theta G}\) away
from the dynamics frame} (\(G\) a random Hermitian of unit spectral
norm; \(\lambda=0.3\); 100 models per row). The columns give the
fraction of models in which records qualify anywhere, and in which the
record-optimal set contains the dynamics frame, the preparation frame,
or neither.

\begin{longtable}[]{@{}
  >{\raggedright\arraybackslash}p{(\columnwidth - 8\tabcolsep) * \real{0.2000}}
  >{\raggedright\arraybackslash}p{(\columnwidth - 8\tabcolsep) * \real{0.2000}}
  >{\raggedright\arraybackslash}p{(\columnwidth - 8\tabcolsep) * \real{0.2000}}
  >{\raggedright\arraybackslash}p{(\columnwidth - 8\tabcolsep) * \real{0.2000}}
  >{\raggedright\arraybackslash}p{(\columnwidth - 8\tabcolsep) * \real{0.2000}}@{}}
\toprule\noalign{}
\begin{minipage}[b]{\linewidth}\raggedright
\(\theta\)
\end{minipage} & \begin{minipage}[b]{\linewidth}\raggedright
Records anywhere
\end{minipage} & \begin{minipage}[b]{\linewidth}\raggedright
In dynamics frame
\end{minipage} & \begin{minipage}[b]{\linewidth}\raggedright
In preparation frame
\end{minipage} & \begin{minipage}[b]{\linewidth}\raggedright
Elsewhere
\end{minipage} \\
\midrule\noalign{}
\endhead
\bottomrule\noalign{}
\endlastfoot
0.0 & 1.00 & 1.00 & 1.00 & 0.00 \\
0.1 & 1.00 & 1.00 & 1.00 & 0.00 \\
0.2 & 1.00 & 1.00 & 0.97 & 0.00 \\
0.4 & 1.00 & 1.00 & 0.44 & 0.00 \\
0.8 & 0.86 & 0.82 & 0.06 & 0.00 \\
1.6 & 0.13 & 0.10 & 0.00 & 0.03 \\
\end{longtable}

At \(\theta=0\) the ``preparation frame'' candidate equals the identity,
a duplicated class of exactly the kind Lemma 3 requires to be merged. As
the ready resource is rotated away, records do not follow it. They stay
in the frame of the coupling and then disappear, in this specific
exploratory family. Finite-angle exponentials of normalised random
Hermitians are not Haar draws, and no Haar limit for this
parameterisation has been established.

\Needspace{12\baselineskip}

\textbf{Table S6. Coupling and ready resource moved together}
(\(K_{\mathrm C}\mapsto U_cK_{\mathrm C}U_c^\dagger\), initial state
\(U_c\left| 000\right\rangle\), \(U_c=e^{-i\varphi G}\); 60 models per
row). The internal-locality minimum was at the identity in every model.

\begin{longtable}[]{@{}
  >{\raggedright\arraybackslash}p{(\columnwidth - 10\tabcolsep) * \real{0.1667}}
  >{\raggedright\arraybackslash}p{(\columnwidth - 10\tabcolsep) * \real{0.1667}}
  >{\raggedright\arraybackslash}p{(\columnwidth - 10\tabcolsep) * \real{0.1667}}
  >{\raggedright\arraybackslash}p{(\columnwidth - 10\tabcolsep) * \real{0.1667}}
  >{\raggedright\arraybackslash}p{(\columnwidth - 10\tabcolsep) * \real{0.1667}}
  >{\raggedright\arraybackslash}p{(\columnwidth - 10\tabcolsep) * \real{0.1667}}@{}}
\toprule\noalign{}
\begin{minipage}[b]{\linewidth}\raggedright
\(\lambda\)
\end{minipage} & \begin{minipage}[b]{\linewidth}\raggedright
\(\varphi\)
\end{minipage} & \begin{minipage}[b]{\linewidth}\raggedright
Records anywhere
\end{minipage} & \begin{minipage}[b]{\linewidth}\raggedright
In internal-locality frame
\end{minipage} & \begin{minipage}[b]{\linewidth}\raggedright
In coupling frame
\end{minipage} & \begin{minipage}[b]{\linewidth}\raggedright
Elsewhere
\end{minipage} \\
\midrule\noalign{}
\endhead
\bottomrule\noalign{}
\endlastfoot
0.3 & 0.4 & 1.000 & 0.583 & 1.000 & 0.000 \\
0.3 & 0.8 & 1.000 & 0.017 & 1.000 & 0.000 \\
1.0 & 0.4 & 1.000 & 0.333 & 0.917 & 0.000 \\
1.0 & 0.8 & 1.000 & 0.033 & 1.000 & 0.000 \\
4.0 & 0.4 & 0.317 & 0.133 & 0.183 & 0.067 \\
4.0 & 0.8 & 0.200 & 0.033 & 0.100 & 0.067 \\
\end{longtable}

The total record-generating Hamiltonian \(H_x\) was more local, by \(L\)
averaged over \(x\), in the coupling frame than in the internal-locality
frame in 100\%, 100\% and 95\% of models at \(\lambda=0.3\), \(1.0\) and
\(4.0\) (\(\varphi=0.8\), 60 models each).

\subsubsection{S2.5 Fresh-seed repeats of the frozen exploratory
configuration}\label{s2.5-fresh-seed-repeats-of-the-frozen-exploratory-configuration}

With the configuration frozen, the main contrasts were re-run with seeds
not used during development.

\begin{itemize}
\tightlist
\item
  \textbf{Seeds 20260926 and 20260927, \(\lambda=0.3\), 200 models per
  arm.} The shared arm again gave \(\bar A=1.000\), \(\bar q=0.050\) and
  \(\hat d=0.950\). The decoupled arm gave \(\hat d=0.072\)
  \([-0.120,0.264]\) and \(0.040\) \([-0.153,0.232]\).
\item
  \textbf{Three-frame design, \(\lambda=0.3\), \(\varphi=0.8\), 100
  models per seed.} Records qualified in every model. They lay in the
  coupling frame in every model and in the internal-locality frame in
  0\% and 1\% of models.
\item
  \textbf{Total Hamiltonian.} \(H_x\) was more local in the coupling
  frame in 100\% of models for both seeds.
\item
  \textbf{Larger model-sample decoupled run, still three qubits.} One
  further run was declared before execution (seed 20260928, 2,952
  models, chosen for a Hoeffding half-width of \(0.05\)). It gave
  records in 10.7\% of models, \(\bar A=0.0515\), \(\bar q=0.0055\) and
  \(\hat d=0.046\) \([-0.004,0.096]\).
\end{itemize}

\subsubsection{S2.6 What the calculation shows, and what it does
not}\label{s2.6-what-the-calculation-shows-and-what-it-does-not}

\begin{enumerate}
\def\labelenumi{\arabic{enumi}.}
\tightlist
\item
  \textbf{Shared frames produce enrichment by construction.} The
  estimate \(\hat d=0.95\) reflects the design, not a regularity of the
  model class.
\item
  \textbf{Records track the frame in which the ready resource and the
  source coupling are aligned.} A ready resource rotated away from the
  coupling produces few records, and they do not follow the resource
  (Table S5).
\item
  \textbf{The verdict depends on which Hamiltonian defines locality.}
  When the coupling and the resource move together, records move with
  them, away from the internal-locality frame, and into the frame in
  which the total Hamiltonian is most local (Table S6 and Section S2.5).
  Its descriptive pattern differs for internal versus branchwise total
  locality. Neither reading is a calibrated verdict on LRA-int or
  LRA-tot, because the exploratory time-averaged record outcome and
  label reference are not the prospective persistent-record test.
\item
  \textbf{The ``decoupled'' arm is only partly decoupled.} Its coupling
  remains local in the dynamics frame. Its small descriptive enrichment
  (\(\hat d=0.046\)) is compatible with influence from the shared
  coupling frame, but the toy does not isolate that cause; and its
  interval, \(U=0.096\), lies just below a margin of \(\Delta=0.1\).
  This is not evidence about any conjecture.
\end{enumerate}

The calculation has three qubits, 20 candidates, Haar candidates not
demonstrated to form an exhaustive set of equivalence classes, a
simplified record score and thresholds chosen during development. The
coupling has a larger declared norm than the internal Hamiltonian at
\(\lambda\le1\); that norm comparison alone is not a record-persistence
guarantee. Its role is to show that the questions in Section 7 need to
be posed with frames declared separately.

\subsection{S3. Post hoc sign tests and window
diagnostics}\label{s3.-post-hoc-sign-tests-and-window-diagnostics}

The four-row directional analysis uses the immutable per-unit contrast
data and verifies them against the full trajectories. Its null is
Pr(z\textgreater0 given z is nonzero)=1/2 for independent signs; it is
not a zero-mean null. No ties were observed. Exact binomial p-values and
a four-row Holm adjustment are reported in the results log. The Holm
family was chosen during review, not in the pre-execution plan.
Dependence between the paired arms is not treated as extra sample size.

The script also reports initial-eligibility and transient-access counts
by frame, and compares both source encodings in an ideal exchange-chain
transport control. The phase-superposition pilot code has singleton
distinguishability \textbar f\textbar, while the
vacuum-versus-excitation control has \textbar f\textbar{} squared. These
are not interchangeable time-window diagnostics. The original study's
thresholds, windows and data remain unchanged.

\begin{Shaded}
\begin{Highlighting}[]
\NormalTok{OPENBLAS\_NUM\_THREADS=1 OMP\_NUM\_THREADS=1 python code/analyze\_review\_additions.py {-}{-}root . {-}{-}out /path/to/fresh\_review\_analysis}
\end{Highlighting}
\end{Shaded}

The code refuses a nonempty output folder. It writes
\nolinkurl{review_analysis.json}, \nolinkurl{sign_tests.csv},
\nolinkurl{transport_window_checks.csv},
\nolinkurl{transport_encodings.csv} and
\nolinkurl{initial_floor_diagnostics.csv}. The last file counts
site/unit pairs, not independent observations. Mathematical derivations,
same-code replay, separate-method checks and external replication remain
distinct.

\subsubsection{S3.1 Initial eligibility and transient strong
access}\label{s3.1-initial-eligibility-and-transient-strong-access}

\begin{longtable}[]{@{}
  >{\raggedright\arraybackslash}p{(\columnwidth - 8\tabcolsep) * \real{0.2000}}
  >{\raggedright\arraybackslash}p{(\columnwidth - 8\tabcolsep) * \real{0.2000}}
  >{\raggedright\arraybackslash}p{(\columnwidth - 8\tabcolsep) * \real{0.2000}}
  >{\raggedright\arraybackslash}p{(\columnwidth - 8\tabcolsep) * \real{0.2000}}
  >{\raggedright\arraybackslash}p{(\columnwidth - 8\tabcolsep) * \real{0.2000}}@{}}
\toprule\noalign{}
\begin{minipage}[b]{\linewidth}\raggedright
Qubits
\end{minipage} & \begin{minipage}[b]{\linewidth}\raggedright
Arm
\end{minipage} & \begin{minipage}[b]{\linewidth}\raggedright
Frame
\end{minipage} & \begin{minipage}[b]{\linewidth}\raggedright
Eligible / total site-unit pairs
\end{minipage} & \begin{minipage}[b]{\linewidth}\raggedright
Ever reaches d \textgreater= 0.8
\end{minipage} \\
\midrule\noalign{}
\endhead
\bottomrule\noalign{}
\endlastfoot
4 & native & L & 288 / 384 & 5 \\
4 & native & C & 192 / 384 & 0 \\
4 & rotated & L & 96 / 384 & 0 \\
4 & rotated & C & 288 / 384 & 0 \\
6 & native & L & 480 / 576 & 10 \\
6 & native & C & 384 / 576 & 0 \\
6 & rotated & L & 288 / 576 & 0 \\
6 & rotated & C & 480 / 576 & 0 \\
\end{longtable}

A site-unit pair is not an independent model. These are post hoc
diagnostic counts from existing trajectories. A transient threshold
crossing does not establish full-window persistence or redundant
copying. The strict persistent-record result remains zero under the
original definition.

\subsection{S4. Factorial extension and paired
uncertainty}\label{s4.-factorial-extension-and-paired-uncertainty}

\subsubsection{S4.1 Design and replay}\label{s4.1-design-and-replay}

The review extension crosses injection L/C with Hamiltonian orientation
L/C, with \(H_C=U_C H_L U_C^\dagger\). Every unit supplies all four
outcomes. The time grid is 0.05 through 2.5 plus the switch-off readout;
the common injection code and two-layer circuit follow the original
source-off conventions. The main paper and results log define z and
\(e_a\). We retain supplied inputs in \nolinkurl{feedback/current/}
without editing their scripts or their predictions.

\nolinkurl{code/replay_reviewer.py} runs the supplied functions with
their original seeds and sample sizes and retains per-unit observables,
rather than only aggregate summaries. It replays the 41,000 factorial
units, the 60,000-unit reimplementation and its smaller checks, and the
6,000-unit third-frame comparison. Each arm in a unit remains paired.
The statement that two routines were written independently refers to
source implementation, not statistical independence of their authors'
reasoning or external peer review.

\subsubsection{S4.2 Range and finite-family
bound}\label{s4.2-range-and-finite-family-bound}

At fixed injection, the initial gap cancels in the difference between
dynamics arms:

\[e_a=\tfrac12\bigl(g_L^a-g_C^a\bigr)\in[-1,1].\]

Let \(s_a^2\) be the unbiased sample variance of e in N iid units.
Theorem 4 of Maurer and Pontil (2009) applies to variables in {[}0,1{]}.
Rescale by \(Y=(e+1)/2\), use both tails, and allocate total error
\(\alpha\) over K specified targets. A simultaneous interval for each
mean is its sample mean plus/minus

\[r_N=\sqrt{\frac{2s_a^2\log(4K/\alpha)}N}+\frac{14\log(4K/\alpha)}{3(N-1)}.\]

\textbf{Derivation.} The theorem's one-tail bound is
\(\sqrt{2V_N\log(2/\delta)/N}+7\log(2/\delta)/(3(N-1))\). Set
\(\delta=\alpha/(2K)\), note \(V_N(Y)=s_a^2/4\), multiply by two to
return to e, and union over the 2K events. No independence between
targets is required. The iid assumption is within each model-unit
sample; paired arms are used to form one e. The variance identity
\(V_N=N^{-1}(N-1)^{-1}\sum_{i<j}(Y_i-Y_j)^2\) equals the usual unbiased
variance. Numerical rounding is not covered by the sampling theorem.

For K=6 and alpha=0.05, all four intervals at N=20,000 exclude zero. The
two eight-qubit intervals at N=1,000 do not. Full normal, Hoeffding and
Bernstein intervals are in
\nolinkurl{data/factorial_analysis/factorial_intervals.csv}. The method
was chosen during review; this is not an external preregistration or a
claim that arbitrary post hoc selections retain coverage.

\subsubsection{S4.3 Important interpretation
checks}\label{s4.3-important-interpretation-checks}

The original negative change relative to frozen dynamics and the
positive Hamiltonian-intervention contrast concern different
counterfactuals. Neither invalidates the other. Isospectrality preserves
the energy eigenvalues, not the initial energy distribution or all
resources relative to the fixed input. The third-frame comparator is an
additional sampled intervention, not a certified link-only control. A
Lieb-Robinson bound is not a sign theorem for e. Regression on an
initial gap must account for the gap's explicit minus-one-half
contribution to z.

\textbf{Reproduction commands, from the extracted release root}

\begin{Shaded}
\begin{Highlighting}[]
\NormalTok{OPENBLAS\_NUM\_THREADS=1 OMP\_NUM\_THREADS=1 python code/replay\_reviewer.py {-}{-}root . {-}{-}out /path/to/fresh\_reviewer\_replay {-}{-}task all}
\NormalTok{python code/analyze\_factorial.py {-}{-}root . {-}{-}out /path/to/fresh\_factorial\_analysis}
\NormalTok{OPENBLAS\_NUM\_THREADS=1 python code/validate\_factorial.py {-}{-}root . {-}{-}out /path/to/fresh\_factorial\_checks.json}
\end{Highlighting}
\end{Shaded}

The analysis reads the delivered \nolinkurl{data/reviewer_replay} by
default. To analyse a replay, use a separate copy with the fresh replay
at that relative location. Do not overwrite the delivered evidence. The
README records exact environment and tool versions; portability to an
arbitrary numerical environment is not assumed. All original scripts can
also be run directly in a copied input folder, with the single-thread
setting stated by their author.

\subsection{S5. Historical execution and
provenance}\label{s5.-historical-execution-and-provenance}

The following subsections preserve the earlier account, with their
historical dates and sample sizes. These are not fresh counts of
independent physical experiments.

\subsubsection{S5.1 Scope}\label{s5.1-scope}

This is supporting material for \emph{Frame-Aligned Records}, review
manuscript v0.5, not an independent research paper. It records new
executions, supplied-script reproduction, prior reported results and
what was not done. The original LRA population study remains unexecuted.
A separately defined source-off feasibility pilot, P01-ID1, has now run;
see the dedicated results log and protocol. It is not silently
relabelled as the original study. No participant data, cosmological
likelihood fit or unrestricted subsystem search was analysed. The
mathematical propositions are assessed through their proofs; numerical
checks are finite safeguards, not universal demonstrations.

\Needspace{12\baselineskip}

\textbf{Table S1. Execution levels must remain separate}

\begin{longtable}[]{@{}
  >{\raggedright\arraybackslash}p{(\columnwidth - 4\tabcolsep) * \real{0.3333}}
  >{\raggedright\arraybackslash}p{(\columnwidth - 4\tabcolsep) * \real{0.3333}}
  >{\raggedright\arraybackslash}p{(\columnwidth - 4\tabcolsep) * \real{0.3333}}@{}}
\toprule\noalign{}
\begin{minipage}[b]{\linewidth}\raggedright
Layer
\end{minipage} & \begin{minipage}[b]{\linewidth}\raggedright
Work done in v1.6
\end{minipage} & \begin{minipage}[b]{\linewidth}\raggedright
Limit
\end{minipage} \\
\midrule\noalign{}
\endhead
\bottomrule\noalign{}
\endlastfoot
Supplied post-v1.5 Claude audit & Executed both numerical scripts in a
fresh directory and compared their two JSON outputs & Same-code replay,
not a new external implementation \\
New CAL-1.6 suite & 174 assertions, including random-case checks and
exact small examples & Related assertions, not 174 independent
experiments \\
Separate CAL-1.6 checker & 46 checks using independent small formulas
and enumeration; does not import the main script & Same-project second
implementation, not external scientific replication \\
Prior v1.5 results & Preserved unchanged with source files and
historical reports & Its large historical assertion counts are not fresh
v1.6 executions \\
Proposed LRA-count experiment & Specification and uncertainty handling
revised & Not run, registered or scientifically confirmed \\
\end{longtable}

\subsubsection{S5.2 Replay of the supplied
audit}\label{s5.2-replay-of-the-supplied-audit}

The exact delivered audit scripts are retained under
\nolinkurl{historical/supporting_lineage/v1_6_code/supplied_claude_audit/}.
Both ran successfully from an initially empty output directory.
\nolinkurl{v15_independent_checks.py} took about 43 seconds and
\nolinkurl{coupling_dominance.py} about 2 seconds in this environment.
Runtime is descriptive, not a portable performance promise. Captured
output and environment differences are recorded in
\nolinkurl{historical/supporting_lineage/v1_6_verification/claude_replay/}.

Recursive comparisons of dictionaries, lists and numerical leaves found
agreement with both supplied JSON outputs using relative tolerance
\(10^{-8}\) and absolute tolerance \(10^{-10}\). Key sets and list
lengths must agree; timing and unrelated files are not used to infer
numerical correctness. The comparison report is
\nolinkurl{historical/supporting_lineage/v1_6_verification/claude_replay_comparison.json}.

The replay includes the supplied exact four-label counterexample, mixed
twirl, gauge invariances, local-persistence example, symbolic examples
and the reported permutation Monte Carlo. For 30 units, the displayed
derangement rejection estimate is approximately 0.060 at a nominal 0.05
threshold; that is a finite Monte Carlo estimate with sampling
uncertainty, not an exact universal size. The exact four-label example
separately establishes a failure of the proposed general validity claim.

The 9,000 perturbation checks all satisfy their bound within the script
tolerance. Those are fragment-time checks across related models, not
9,000 independent experiments. They confirm the stated implementation of
the inequality, not a theorem that coupling-frame records must maximize
a locality-record score.

\subsubsection{S5.3 New suite and effective
configuration}\label{s5.3-new-suite-and-effective-configuration}

The new executable source is
\nolinkurl{historical/supporting_lineage/v1_6_code/checks_v0_3_0.py}. It
accepts one complete JSON configuration, rejects missing or unused keys,
checks numerical ranges, refuses a nonempty output directory, and writes
the effective configuration and its canonical-JSON hash with every run.
The version is CAL-1.6; the code version is 0.3.0. This is calibration
code, not the future LRA experiment implementation.

\Needspace{12\baselineskip}

\textbf{Table S2. Complete new calibration configuration}

\begin{longtable}[]{@{}
  >{\raggedright\arraybackslash}p{(\columnwidth - 4\tabcolsep) * \real{0.3333}}
  >{\raggedright\arraybackslash}p{(\columnwidth - 4\tabcolsep) * \real{0.3333}}
  >{\raggedright\arraybackslash}p{(\columnwidth - 4\tabcolsep) * \real{0.3333}}@{}}
\toprule\noalign{}
\begin{minipage}[b]{\linewidth}\raggedright
Field
\end{minipage} & \begin{minipage}[b]{\linewidth}\raggedright
Value
\end{minipage} & \begin{minipage}[b]{\linewidth}\raggedright
Use
\end{minipage} \\
\midrule\noalign{}
\endhead
\bottomrule\noalign{}
\endlastfoot
version & CAL-1.6 & Validate scope \\
seed & 20260926 & NumPy default RNG stream \\
random\_cases & 32 & Each random matrix/invariance/perturbation check
family \\
mixed\_draws & 1200 & Fixed-spectrum twirl Monte Carlo \\
word\_length & 5 & Symbolic support enumeration \\
chain\_length & 30000 & Finite golden-mean example \\
bound\_epsilon & 0.1 & Adaptive-family probability envelope \\
net\_radius & 0.05 & Operator-norm cover radius \\
angle\_limit & pi & Declared circuit parameter interval \\
max\_bound\_qubits & 32 & Envelope evaluation through 32 qubits; no
32-qubit state simulation \\
\end{longtable}

The canonical effective-configuration SHA-256 is
\nolinkurl{f5f8d1633e36e402922eefb5fe8a24640b890fa028e0985856a753f779cabc2b}.
Numerical environment: Python 3.13.5, NumPy 2.3.5 and SciPy 1.17.0. The
supplied Claude audit originally reports Python 3.11.15, NumPy 2.4.4 and
SciPy 1.17.1, so agreement is assessed with tolerances, not assumed from
version identity. The code's environment report supplies the actual
current values rather than promising installation compatibility with
every platform.

\Needspace{12\baselineskip}

\textbf{Table S3. Data files and their unit of observation}

\begin{longtable}[]{@{}
  >{\raggedright\arraybackslash}p{(\columnwidth - 4\tabcolsep) * \real{0.3333}}
  >{\raggedright\arraybackslash}p{(\columnwidth - 4\tabcolsep) * \real{0.3333}}
  >{\raggedright\arraybackslash}p{(\columnwidth - 4\tabcolsep) * \real{0.3333}}@{}}
\toprule\noalign{}
\begin{minipage}[b]{\linewidth}\raggedright
File in historical/supporting\_lineage/v1\_6\_data/
\end{minipage} & \begin{minipage}[b]{\linewidth}\raggedright
Contents
\end{minipage} & \begin{minipage}[b]{\linewidth}\raggedright
Interpretation
\end{minipage} \\
\midrule\noalign{}
\endhead
\bottomrule\noalign{}
\endlastfoot
per\_case\_checks.csv & 32 cases, gauge and stabilizer residuals, local
bound and sampled slope & Numerical checks of deterministic
inequalities \\
perturbation\_checks.csv & 32 differences and perturbation bounds & New
small finite checks, distinct from the replayed 9,000 checks \\
localized\_constants.csv & Global and endpoint touching constants for 2,
4, 6 qubits & One explicitly specified bounded-strength model \\
mixed\_draws.csv & 1,200 squared Hilbert-Schmidt deviations & Individual
draws, not only an aggregate \\
symbolic\_support.csv & All 32 binary length-five words and their model
probabilities & Exact support and finite-sample occurrence are distinct
columns \\
adaptive\_cover\_bound.csv & Log probability envelopes for n=4 through
32 & Evaluation of a theorem bound, not physical observations \\
results.json & Assertions, selected summaries, environment and code hash
& Failures cannot be hidden by counting only passed examples \\
\end{longtable}

The old constant-window example retains its analytic certificate
0.0097751350, before adding a numerical error budget. The new endpoint
chain has \(K_F=\sqrt{1.0025}\) for every plotted size. The supplied
random-chain example is a different construction, not a failed attempt
to reproduce this value.

\subsubsection{S5.4 Independent small verification and reproduction
instructions}\label{s5.4-independent-small-verification-and-reproduction-instructions}

The second checker does not import the main checker. It reconstructs
exact permutation ranks, contracts the twirl in a different form,
evaluates endpoint spectra analytically, and checks the circuit-cover
arithmetic. Its 46 checks are another implementation path within the
same session, not a second research team.

From the extracted release root, use initially absent output
directories:

\begin{Shaded}
\begin{Highlighting}[]
\NormalTok{python historical/supporting\_lineage/v1\_6\_code/checks\_v0\_3\_0.py {-}{-}config historical/supporting\_lineage/v1\_6\_code/config.json {-}{-}output /path/to/fresh\_calibration}
\NormalTok{python historical/supporting\_lineage/v1\_6\_code/second\_checks\_v0\_3\_0.py {-}{-}results /path/to/fresh\_calibration}
\NormalTok{python historical/supporting\_lineage/v1\_6\_code/plot\_results.py {-}{-}data /path/to/fresh\_calibration {-}{-}output /path/to/fresh\_figures}
\end{Highlighting}
\end{Shaded}

The legacy scripts are preserved for provenance and old calibration
replay. They are not silently relabelled as corrected production
software. In particular, the earlier hard-coded defaults and
aggregate-only outputs remain historical limitations. New matrix cases
and mixed draws have individual rows; a future model-search
implementation must additionally log every input frame, rejection, tie,
search failure and certificate interval.

For raw floating-point checks, a small residual is not a verified
interval enclosure. The many-body continuum certificate requires a
defensible numerical allowance or interval method. Without one, report
numerical evidence and a certificate conditional on the allowance, not
rigorous machine certification.

\subsubsection{S5.5 Conditional pairing retained only as an
alternative}\label{s5.5-conditional-pairing-retained-only-as-an-alternative}

The v1.6 primary prospective target compares normalized counts within an
input-eligible model. It uses no cross-model pairing. The previous
pairing method remains meaningful only under its own contract: specify
the independent unit, frame registration, conditioning variables,
admissible group action, why the null is invariant, tie handling and
missing-output treatment. A complete permutation group cannot create
exchangeability in a physical population that lacks it.

Fixed common frames can make matched and mismatched distances identical.
Randomly varying shared orientations can produce an apparent pairing
signal inherited from supplied structure. These are valid diagnostic
examples, not reasons to ban every permutation test. Where the null
cannot be justified, the appropriate status is
NO\_IDENTIFIED\_PAIRING\_NULL. The original conditional wording is
preserved in \nolinkurl{editorial/retained_pairing_contract.md}.

\subsubsection{S5.6 Delivery and scientific
readiness}\label{s5.6-delivery-and-scientific-readiness}

The present supplement is complete for inspection at its declared finite
scope. The full release contains the exact new code and data plus the
supplied audit and baseline archive. A short manuscript-only review
necessarily cannot certify those absent files, and no claim in the text
turns that limitation into disconfirmation. Source parity, render
inspection, archive integrity and scientific correctness are separate
checks, recorded in the release verification report.

The new uniform-cover theorem is a conditional mathematical corollary of
known concentration and covering arguments. Its novelty and physical
relevance remain open. The broad AIU inquiry neither inherits empirical
confirmation from it nor is refuted by a poor finite-size bound.

\subsubsection{S5.7 Maintenance patch and separate source-off
execution}\label{s5.7-maintenance-patch-and-separate-source-off-execution}

Appendix C of the methods paper retains the archived v1.5 2,000-draw
twirl result (mean 0.0557464152; standard error 0.0007264476). CAL-1.6
used 1,200 draws instead (mean 0.05550839355; standard error
0.00093479233), with exact target 0.056. These are two calibration
samples, not a conflict or one enlarged study. The historical data paths
identify them separately.

P01-ID1 is a new finite computational pilot with 96 independent model
units at each of four and six sites. Each unit supplies two paired
injection arms and paired evolution-on/frozen controls. Its primary
outcome is a baseline-subtracted checkpoint-average local-readability
contrast, not persistent redundancy, the older LRA-count target, or a
recovered TPS. Full trajectories, model coefficients, random-stream
identifiers, configuration and code hashes are in
\nolinkurl{data/pilot_run/}; the specification and local pre-execution
receipt are in \nolinkurl{pilot/}. The specialist review was not sent
and no external validation is claimed.

The main solver diagonalizes each internal Hamiltonian.
\nolinkurl{code/validate_injection_pilot.py} independently checks
selected trajectories with exponential action and a full-density-matrix
partial trace. It also checks exact transport, source-free global
distinguishability, zero evolution, passive conjugation and closed
stationary records. Its finite-precision agreement is not a rigorous
roundoff enclosure. Persistence classifications are conditional on the
explicitly empirical numerical guard; failed certificates are not absent
records.

Reproduction of the new pilot from an empty directory:

\begin{Shaded}
\begin{Highlighting}[]
\NormalTok{OPENBLAS\_NUM\_THREADS=1 OMP\_NUM\_THREADS=1 python code/run\_injection\_pilot.py {-}{-}config pilot/config.json {-}{-}out /path/to/fresh\_run}
\NormalTok{python code/validate\_injection\_pilot.py {-}{-}root /path/to/extracted\_patch}
\end{Highlighting}
\end{Shaded}

The validator reads the delivered \nolinkurl{data/pilot_run} by default.
To validate a replay, place that replay at the same relative path in a
separate extracted copy. Do not overwrite the delivered evidence. The
code and effective configuration specify all model parameters; no unused
fields are silently accepted.

\subsection{References}\label{references}

Maurer, A., and Pontil, M. (2009). Empirical Bernstein Bounds and Sample
Variance Penalization. COLT 2009. \url{https://arxiv.org/abs/0907.3740}

\end{document}
